% Copyright (c) 2026 Geovana Neves. Licensed under CC BY 4.0 (https://creativecommons.org/licenses/by/4.0/). See LICENSE-AND-AUTHORSHIP.md.
%
% 01-fts-overview/text/03-workflow.tex
%
% Part 3: the simulation workflow as ten stages, from the question to the
% engineering analysis, with what each stage decides. Manual pages are those
% of build 26121 [6].

\section{The simulation workflow}

\begin{frame}{Ten stages, in the order decisions constrain each other}
\relax
{\ftstablesize\renewcommand{\arraystretch}{1.12}
\begin{tabularx}{\linewidth}{@{}r@{\hspace{2mm}}P{0.34\linewidth}R@{}}
\toprule
\textbf{\#} & \textbf{Stage} & \textbf{The decision it owns} \\
\midrule
1  & Case study and scope & what question the case answers, and what it does not \\
2  & Matrix definition & which points are run, and why those \\
3  & Water-tight geometry & blunt or sharp trailing edge, and where the components split \\
4  & Surface mesh & where resolution is spent \\
5  & Boundary conditions & which edges shed vorticity, and how much \\
6  & Solver setup & what the solver is allowed to assume \\
7  & Execution and monitoring & whether the answer has converged \\
8  & Output definition and export & what evidence the run leaves behind \\
9  & Verification and validation & whether the answer is right \\
10 & Engineering analyses & what the answer means for the design \\
\bottomrule
\end{tabularx}}
\vspace{1mm}
\begin{takeaway}The stages are not a checklist to tick: a decision taken out of
order has to be taken again.\end{takeaway}
\end{frame}

\begin{frame}{One session in the interface, and the lock in it}
\slidefig[0.34]{g01-workflow-chart}
\begin{itemize}\small
  \item \textbf{Initialisation is a lock.} It freezes the solver against the
        mesh and the boundary conditions as they stand, and fixes the solver
        mode and the symmetry claim [6, pp.~218-221].
  \item A runtime setting (angle, velocity, iterations) can change after it.
        A mesh or boundary change needs the initialisation removed and taken
        again. That is why boundary conditions precede setup [6, p.~411].
\end{itemize}
\end{frame}

\begin{frame}{Stages 3 and 4: geometry and surface mesh}
\begin{columns}[T]
\begin{column}{0.48\textwidth}
\begin{block}{Geometry}
\begin{itemize}\small
  \item \textbf{Water-tight} means the body closes, not the domain.
  \item Collapsed or blunt trailing edge: decided in CAD (Part 2).
  \item Components are split by the \textbf{boundary condition they will
        receive}, and by the loads you will want per component, never by CAD
        convenience.
\end{itemize}
\end{block}
\end{column}
\begin{column}{0.48\textwidth}
\begin{block}{Surface mesh}
\begin{itemize}\small
  \item Four sources: the tool's cross-section meshers, transfinite
        interpolation, the wrapper, or an imported mesh [6, pp.~72-75].
        Only the first two regenerate in seconds.
  \item Faces go where the gradients are: leading edge, tip, junction, the
        row the wake leaves from [6, p.~62].
  \item A mesh convergence study is an obligation, and it is affordable.
\end{itemize}
\end{block}
\end{column}
\end{columns}
\begin{takeaway}The mesh carries the boundary conditions. A boundary is a
decision about what you will be able to ask for later.\end{takeaway}
\end{frame}

\begin{frame}{Stage 5: boundary conditions are marking}
\begin{columns}[T]
\begin{column}{0.50\textwidth}
\begin{itemize}\small
  \item \textbf{Farfield}: there is none to set. \textbf{Walls}: every face is
        a slip wall once meshed, switched to no slip by the solver when the
        boundary layer is active [6, p.~164].
  \item What remains is \textbf{marking} where that default is wrong:
        trailing edges, with their type (Part 2); base regions; wake
        termination nodes; and, when the case needs them, actuator discs,
        inlets and outlets, custom free-stream fields [6, pp.~164-197].
\end{itemize}
\end{column}
\begin{column}{0.46\textwidth}
\begin{alertblock}{Auto-detection is a proposal}
\small The solver can detect trailing edges; it proposes, you decide. Inspect
every marked edge before initialisation: a split trailing edge is two edges,
and a missing one is silent.
\end{alertblock}
\end{column}
\end{columns}
\end{frame}

\begin{frame}{Stage 6: what the solver may assume}
\relax
{\ftstablesize\renewcommand{\arraystretch}{1.12}
\begin{tabularx}{\linewidth}{@{}P{0.14\linewidth}R@{}}
\toprule
\textbf{Symmetry} & \textbf{What it asserts} \\
\midrule
None & nothing: the mesh in front of the solver is the whole problem \\
Mirror & a symmetry plane in the flow, not only in the geometry; false as
  soon as the sideslip is not zero \\
Periodic & an onset flow the same at every azimuth, plus a copy count; false
  as soon as the angle of attack is not zero \\
\bottomrule
\end{tabularx}}
\vspace{1mm}
\begin{itemize}\small
  \item \textbf{Coordinate systems first}: loads are reported in one, and a
        moment without its frame and point is not a number.
  \item \textbf{Reference velocity, area and length} normalise every
        coefficient. Set them; never inherit them.
  \item \textbf{Convergence is three controls}: an iteration cap, a residual
        threshold, and how long the residual must stay below it
        [6, pp.~202-203]. Viscous settings change the answer; leaving them off
        is a decision too.
\end{itemize}
\end{frame}

\begin{frame}{Stages 7 and 8: run, watch, and declare the outputs}
\begin{columns}[T]
\begin{column}{0.48\textwidth}
\begin{block}{Execution and monitoring}
\begin{itemize}\small
  \item \textbf{Converged is a state that was held}, not a number that was
        touched. A run stopped at its cap ran out of iterations.
  \item Unsteady asks twice: the residual inside each step, and the history
        of the quantity you want. They are independent.
  \item Periodic: judge the \textbf{per-revolution} statistic, not the
        sample.
\end{itemize}
\end{block}
\end{column}
\begin{column}{0.48\textwidth}
\begin{block}{Output definition and export}
\begin{itemize}\small
  \item \textbf{An output not declared was not computed}: sections, probes
        and plots are set before the run.
  \item Symmetry-load doubling changes what the solver computes. Two runs set
        differently are not the same quantity.
  \item Every coefficient carries a normalisation: record it with the number.
\end{itemize}
\end{block}
\end{column}
\end{columns}
\begin{takeaway}Recovering a variable that was never declared means running
the case again, which for a rotor is the most expensive mistake
available.\end{takeaway}
\end{frame}

\begin{frame}{Stages 9 and 10: is it right, and what does it mean}
\begin{itemize}\small
  \item \textbf{Verification}: the solution against a theoretical anchor and
        against itself. A finite-wing lift slope, a span efficiency, a mesh
        and a time-step study, a sector against the full wheel.
  \item \textbf{Validation}: the solution against measured data, inside the
        envelope where the method has been compared with experiment
        [13, 14].
  \item An error that scales the whole solution divides out of a ratio. An
        efficiency can agree while thrust and torque are both off: check the
        forces, not only their ratio.
  \item \textbf{Engineering analyses} use the answer: derivatives, trim,
        installation deltas, component breakdowns. They inherit every
        assumption of stages 3 to 8.
\end{itemize}
\begin{takeaway}A mid-fidelity number is defensible when its build, its
settings, its convergence and its verification travel with it.\end{takeaway}
\end{frame}
